\section{Temporal Accuracy and Stability Analysis}
\label{sec:stability_analysis}

\subsection{Semi-Discrete System and Eigenvalue Analysis}

After spatial discretization and implementation of the boundary
conditions, the energy equation can be expressed as a system of
ordinary differential equations:
\begin{equation}
\boxed{
\frac{d\mathbf{T}}{dt}
=
\mathbf{A}\mathbf{T}+\mathbf{b},
}
\label{eq:stability_semidiscrete}
\end{equation}
where $\mathbf{T}$ is the vector of physical-cell temperatures,
$\mathbf{A}$ is the assembled spatial-operator matrix, and
$\mathbf{b}$ contains the temperature-independent contributions
from prescribed boundary values and viscous dissipation.

The matrix $\mathbf{A}$ depends on the spatial discretization,
grid resolution, prescribed velocity field, Peclet number,
convection coefficient, and implementation of the boundary
conditions. Since the computational domain has non-periodic
boundary conditions and a streamwise velocity that varies
across the channel, the eigenvalues cannot generally be
obtained using the Fourier-mode formula for a periodic,
constant-coefficient stencil.

Instead, the spatial matrix must first be assembled using
the actual discretized equations and boundary conditions.
Its eigenvalues are then calculated numerically from
\begin{equation}
\boxed{
\mathbf{A}\mathbf{q}_k
=
\lambda_k\mathbf{q}_k,
\qquad k=1,\ldots,N,
}
\label{eq:spatial_eigenvalue_problem}
\end{equation}
where $\lambda_k$ and $\mathbf{q}_k$ are the eigenvalues
and corresponding eigenvectors, and $N=n_xn_y$ is the
number of physical-cell temperature unknowns.

The eigenvalues characterize the modal behaviour of the
linear spatial operator and provide the basis for examining
the stability of the time-integration methods. Since the
source vector $\mathbf{b}$ is independent of temperature,
it does not directly enter the amplification factors
derived below.

\subsection{Amplification Factors and Temporal Accuracy}

Consider the scalar linear equation
\begin{equation}
\frac{dw}{dt}=\lambda w,
\label{eq:stability_test_equation}
\end{equation}
where $\lambda$ represents an eigenvalue of the spatial
operator. Define the dimensionless parameter
\begin{equation}
z=\lambda\Delta t,
\label{eq:stability_parameter}
\end{equation}
where $\Delta t$ is the time step.

The exact solution advances according to
\begin{equation}
w^{n+1}=e^z w^n.
\label{eq:exact_amplification}
\end{equation}
A numerical time-integration method produces an update of
the form
\begin{equation}
w^{n+1}=G(z)w^n,
\label{eq:numerical_amplification}
\end{equation}
where $G(z)$ is the amplification factor of the method.

Comparing the amplification factor with the Taylor expansion
of the exact solution,
\begin{equation}
e^z
=
1+z+\frac{z^2}{2!}
+\frac{z^3}{3!}
+\frac{z^4}{4!}+\cdots,
\label{eq:exact_taylor_expansion}
\end{equation}
provides a way to determine the temporal order of accuracy.
A method is of order $p$ when its amplification factor
matches the expansion of $e^z$ through the term $z^p$,
with the first unmatched term having order $z^{p+1}$.

\subsubsection{Implicit Euler}

The implicit Euler method evaluates the spatial operator
at the next time level:
\begin{equation}
\frac{w^{n+1}-w^n}{\Delta t}
=
\lambda w^{n+1}.
\end{equation}
Rearranging gives
\begin{equation}
w^{n+1}
=
\frac{1}{1-\lambda\Delta t}w^n.
\end{equation}
Therefore, its amplification factor is
\begin{equation}
\boxed{
G_{\mathrm{IE}}(z)=\frac{1}{1-z}.
}
\label{eq:implicit_euler_amplification}
\end{equation}

Expanding this expression as a power series gives
\begin{equation}
G_{\mathrm{IE}}(z)
=
1+z+z^2+z^3+\cdots.
\end{equation}
Comparison with Equation~\eqref{eq:exact_taylor_expansion}
shows agreement only through the linear term. Implicit
Euler is therefore first-order accurate in time, with a
global temporal error of order $\mathcal{O}(\Delta t)$
for sufficiently smooth solutions.

The absolute-stability condition is
\begin{equation}
\left|
\frac{1}{1-\lambda\Delta t}
\right|
\leq 1.
\label{eq:implicit_euler_stability}
\end{equation}
If $\operatorname{Re}(\lambda)\leq 0$, then
\begin{equation}
\begin{aligned}
|1-\lambda\Delta t|^2
&=
1-2\Delta t\,\operatorname{Re}(\lambda)
+\Delta t^2|\lambda|^2
\\
&\geq 1.
\end{aligned}
\end{equation}
Consequently, implicit Euler satisfies the absolute-stability
condition for every positive time step when the eigenvalues
of the spatial operator lie in the non-positive real
half-plane. This property is known as $A$-stability.

However, $A$-stability does not guarantee that an arbitrarily
large time step will produce an accurate transient solution
or efficiently converge to a steady state. The spectrum of
the actual spatial matrix and the practical convergence
behaviour must still be considered.

\subsubsection{Second-Order Runge--Kutta Method}

The explicit midpoint method advances the solution through
an intermediate state:
\begin{align}
w^{(1)}
&=
w^n+\frac{\Delta t}{2}\lambda w^n,
\\
w^{n+1}
&=
w^n+\Delta t\,\lambda w^{(1)}.
\end{align}
Substituting the intermediate value into the final update
gives
\begin{equation}
w^{n+1}
=
\left(
1+\lambda\Delta t
+\frac{(\lambda\Delta t)^2}{2}
\right)w^n.
\end{equation}
The amplification factor is therefore
\begin{equation}
\boxed{
G_{\mathrm{RK2}}(z)
=
1+z+\frac{z^2}{2}.
}
\label{eq:rk2_amplification}
\end{equation}

This polynomial matches the Taylor expansion of the exact
amplification factor through the quadratic term. The
explicit midpoint method is consequently second-order
accurate in time, with global temporal error
$\mathcal{O}(\Delta t^2)$ for sufficiently smooth solutions.

Its absolute-stability condition is
\begin{equation}
\boxed{
\left|
1+\lambda_k\Delta t
+\frac{(\lambda_k\Delta t)^2}{2}
\right|
\leq 1
\qquad \text{for every } k.
}
\label{eq:rk2_stability}
\end{equation}

Since this condition is not satisfied for arbitrary complex
values of $\lambda_k\Delta t$, the explicit midpoint method
is conditionally stable. The allowable time step depends on
the eigenvalues of the assembled spatial matrix, which in
turn depend on the grid, velocity field, Peclet number,
and convection coefficient.

\subsubsection{Fourth-Order Runge--Kutta Method}

The classical fourth-order Runge--Kutta method evaluates
the right-hand side at four stages:
\begin{align}
k_1&=\lambda w^n,
\\
k_2&=\lambda\left(w^n+\frac{\Delta t}{2}k_1\right),
\\
k_3&=\lambda\left(w^n+\frac{\Delta t}{2}k_2\right),
\\
k_4&=\lambda\left(w^n+\Delta t\,k_3\right).
\end{align}
The solution is updated according to
\begin{equation}
w^{n+1}
=
w^n+\frac{\Delta t}{6}
\left(k_1+2k_2+2k_3+k_4\right).
\end{equation}
Substitution of the four stages yields the amplification
factor
\begin{equation}
\boxed{
G_{\mathrm{RK4}}(z)
=
1+z+\frac{z^2}{2}
+\frac{z^3}{6}
+\frac{z^4}{24}.
}
\label{eq:rk4_amplification}
\end{equation}

This polynomial matches the Taylor expansion of the exact
amplification factor through the fourth-order term.
The classical Runge--Kutta method is therefore fourth-order
accurate in time, with global temporal error
$\mathcal{O}(\Delta t^4)$ for sufficiently smooth solutions.

Its absolute-stability condition is
\begin{equation}
\boxed{
\left|
1+\lambda_k\Delta t
+\frac{(\lambda_k\Delta t)^2}{2}
+\frac{(\lambda_k\Delta t)^3}{6}
+\frac{(\lambda_k\Delta t)^4}{24}
\right|
\leq 1
\qquad \text{for every } k.
}
\label{eq:rk4_stability}
\end{equation}

Like RK2, classical RK4 is conditionally stable. Its
stability region differs from that of RK2, so its allowable
time step must be assessed using its own amplification
factor rather than inferred from the RK2 stability limit.

\subsection{Numerical Assessment of Stability}

Once the spatial matrix has been assembled, the eigenvalues
can be used to assess the modal stability of the explicit
Runge--Kutta methods. For a candidate time step, define
\begin{equation}
G_{\max}(\Delta t)
=
\max_{k=1,\ldots,N}
\left|G(\lambda_k\Delta t)\right|.
\label{eq:maximum_amplification}
\end{equation}
The candidate time step satisfies the modal
absolute-stability criterion if
\begin{equation}
\boxed{
G_{\max}(\Delta t)\leq 1.
}
\label{eq:maximum_amplification_condition}
\end{equation}

The maximum stable time step can be estimated by evaluating
this condition over candidate time steps and refining the
search near the stability boundary. This calculation is
performed separately for RK2 and RK4 because their
amplification factors differ.

The direct sparse solver uses implicit Euler time integration.
Its time integrator is $A$-stable when the spectrum of the
assembled spatial operator lies in the non-positive real
half-plane. The direct linear-system solution itself does not
introduce a separate explicit time-step stability restriction.
Nevertheless, because the solver is used for pseudo-time
marching toward steady state, its practical behaviour depends
on the selected time step, the properties of the discrete
operator, and the numerical conditioning of the linear system.
These factors are assessed through numerical convergence tests.

The approximate factorization method requires a separate
assessment. Although it is derived from the implicit Euler
equation, replacing the combined spatial operator with a
product of directional operators changes the numerical
update. Its stability therefore cannot be inferred solely
from the amplification factor of fully implicit Euler.
In this project, the behaviour of the approximate
factorization method is evaluated numerically by testing
candidate time steps and monitoring the resulting residuals
and convergence behaviour.

The stability analysis therefore provides a
theoretical basis for choosing time steps, while numerical
tests and convergence histories are used to verify the
behaviour of the implemented solvers.